\chapter{Système de décision hybride et confrontation à la notation réelle}

\section{D'un score à une décision}

Les Chapitres~4 et~5 ont produit deux \emph{scores} — une grandeur continue à interpréter. Or
l'objectif opérationnel posé pour ce projet (Chapitre~1) est différent : produire un
\textbf{système de décision}, c'est-à-dire une sortie discrète et actionnable
(accepté / refusé / à étudier), assortie d'une justification explicite. Cette distinction reprend
la terminologie usuelle en aide à la décision : la banque dispose aujourd'hui d'un système
d'\emph{aide} à la décision (un score que l'analyste doit interpréter) ; ce projet propose un
système de \emph{décision} au sens où il rend un verdict directement exploitable, tout en
conservant une zone d'arbitrage humain explicite pour les cas ambigus.

\section{Score hybride pour la notation A-E}

Un premier score hybride, destiné à reconstruire la notation-lettre MNS2 (A à E), combine les
trois sorties disponibles par une moyenne pondérée :
\[
S_{\text{hybride}} = 0{,}30 \cdot S_{\text{expert}} + 0{,}35 \cdot S_{\text{RF}} + 0{,}35 \cdot S_{\text{LogReg}}
\]
où $S_{\text{RF}}$ et $S_{\text{LogReg}}$ sont des scores sur 100 dérivés des probabilités de
classe prédites par les deux modèles (Chapitre~5, adaptés ici à la prédiction multi-classe A-E).
Le Tableau~\ref{tab:comparaison-mns2} confronte, pour chaque dossier, la notation-lettre réelle à
la notation prédite par le système hybride.

\begin{table}[h]
\centering
\small
\begin{tabular}{@{}lcccc@{}}
\toprule
\textbf{Dossier} & \textbf{MNS2 réel} & \textbf{Score hybride /100} & \textbf{Notation prédite} & \textbf{Divergence} \\
\midrule
Feuil2  & E & 52.9 & C & \textcolor{bporange}{\textbf{Oui}} \\
Feuil1  & D & 51.5 & C & Non \\
Feuil11 & B & 50.1 & C & Non \\
Feuil13 & D & 53.8 & C & Non \\
Feuil3  & C & 50.7 & C & Non \\
Feuil12 & C & 44.6 & D & Non \\
Feuil4  & D & 42.9 & D & Non \\
Feuil5  & D & 47.2 & D & Non \\
Feuil8  & E & 42.8 & D & Non \\
\bottomrule
\end{tabular}
\caption{Confrontation notation MNS2 réelle vs. notation prédite (score hybride)}
\label{tab:comparaison-mns2}
\end{table}

Un seul dossier — \textbf{Feuil2} — présente un écart de deux crans ou plus (E réel contre C
prédit), et est à ce titre flaggé automatiquement par le système comme divergence à investiguer.
Les autres dossiers présentent des écarts d'au plus un cran, ce qui, compte tenu de la taille de
l'échantillon, ne permet pas de conclure formellement à une bonne performance du système mais
constitue un signal encourageant sur la cohérence d'ensemble entre logique experte, apprentissage
statistique et notation bancaire réelle.

\begin{limite}[Cas Feuil2 à investiguer]
Le dossier Feuil2 présente une notation MNS2 réelle très défavorable (E) alors que le score
expert (76,4/100, cf. Tableau~4.2) et le système hybride (52,9/100, notation C) le classent de
façon nettement plus favorable. Trois explications sont possibles : (i) un facteur qualitatif non
capturé par les variables quantitatives disponibles a pesé fortement dans la décision bancaire
réelle ; (ii) une erreur de segment ou de saisie affecte ce dossier ; (iii) la pondération de la
grille experte (Chapitre~4) est mal calibrée pour ce profil de dossier. Ce cas est identifié comme
prioritaire pour une investigation qualitative avec le maître de stage.
\end{limite}

\section{Moteur de décision logique}

Pour la détection de défaut (Chapitre~5), le moteur de décision applique une règle explicite,
combinant la probabilité de défaut hybride et le score expert comme garde-fou :

\begin{definition}{Règle de décision}
Soit $p = \dfrac{p_{\text{RF}} + p_{\text{LogReg}}}{2}$ la probabilité de défaut hybride (moyenne
des deux modèles lorsque les deux sont disponibles). La décision du système est :
\[
\text{Décision} =
\begin{cases}
\text{REFUSÉ} & \text{si } p \geq 0{,}60 \\
\text{ACCEPTÉ} & \text{si } p \leq 0{,}25 \text{ et } S_{\text{expert}} \geq 30 \\
\text{À ÉTUDIER} & \text{sinon}
\end{cases}
\]
Le garde-fou sur $S_{\text{expert}}$ empêche une acceptation automatique lorsque le score
expert est très défavorable malgré un avis favorable du modèle statistique — une divergence entre
les deux approches (Chapitre~2, Section 2.5) est alors traitée comme un signal d'alerte à
soumettre à l'arbitrage humain plutôt que comme un cas tranché.
\end{definition}

Les seuils (60\,\% / 25\,\%) sont, comme les pondérations de la grille experte, des
\textbf{hypothèses de travail} à ajuster avec le maître de stage — par exemple par calibration
sur un échantillon de dossiers plus large, une fois le workflow applicatif (Section~6.4)
alimenté par de nouveaux cas.

\section{Résultats du système de décision}

\begin{table}[h]
\centering
\small
\begin{tabular}{@{}lccc@{}}
\toprule
\textbf{Dossier} & \textbf{P(défaut) hybride} & \textbf{Décision} & \textbf{Cohérence avec label} \\
\midrule
Feuil8  & 88.9\,\% & REFUSÉ & Non (label = sain) \\
Feuil11 & 85.8\,\% & REFUSÉ & Non (label = sain) \\
Feuil13 & 71.6\,\% & REFUSÉ & Oui (label = défaut) \\
Feuil2  & 52.0\,\% & À ÉTUDIER & — \\
Feuil5  & 43.3\,\% & À ÉTUDIER & — \\
Feuil1  & 41.5\,\% & À ÉTUDIER & — \\
Feuil3  & 41.0\,\% & À ÉTUDIER & — \\
Feuil12 & N/A & À ÉTUDIER & Données insuffisantes \\
Feuil4  & N/A & À ÉTUDIER & Données insuffisantes \\
\bottomrule
\end{tabular}
\caption{Décisions rendues par le système}
\label{tab:decisions}
\end{table}

Sur cet échantillon, le système ne rend jamais de décision « ACCEPTÉ » automatique : les
probabilités de défaut estimées, même pour les dossiers labellisés sains, restent au-dessus du
seuil de 25\,\%. Ce constat, plutôt qu'une faiblesse du système, illustre une conséquence directe
et attendue de la taille de l'échantillon : avec seulement 6 observations d'entraînement à
chaque itération LOOCV, les modèles restent prudents et ne produisent pas de probabilités très
proches de 0. Le système privilégie ainsi, par construction, l'arbitrage humain (« à étudier »)
à la décision automatique dans les cas ambigus — un choix de conception délibéré compte tenu des
enjeux d'un refus ou d'un octroi de crédit erroné.

\section{Workflow applicatif et traçabilité}

Le système de décision est intégré à une application (tableau de bord interactif) organisée
autour de trois vues :

\begin{itemize}
    \item \textbf{Analyste} : dépôt d'un dossier et rédaction d'un commentaire argumentant en
    faveur du client (element qualitatif non capturé par les variables quantitatives) ;
    \item \textbf{Direction} : consultation de la décision automatique et de sa justification,
    croisée avec l'argumentaire de l'analyste, puis validation ou arbitrage final ; cette étape
    génère automatiquement un procès-verbal (PV) de décision, horodaté et exportable ;
    \item \textbf{Explicabilité} : vue anonymisée présentant les facteurs déterminants
    (Chapitre~5, Section 5.4) et l'ensemble des décisions rendues, à des fins d'audit et de
    pédagogie interne.
\end{itemize}

Ce workflow répond directement à la problématique de transparence posée au Chapitre~1 : chaque
décision est traçable, justifiée, et laisse une place explicite à l'expertise humaine sur les cas
ambigus, sans reproduire l'opacité de « boîte noire » identifiée dans la littérature
(Chapitre~2, Section 2.4).

\section{Synthèse}

La confrontation du système hybride à la notation MNS2 réelle fait apparaître une cohérence
globale satisfaisante (huit dossiers sur neuf sans divergence marquante), tout en identifiant un
cas prioritaire d'investigation qualitative (Feuil2). Le moteur de décision transforme ces scores
en verdicts actionnables, avec une zone d'arbitrage humain explicite pour les cas où la confiance
du système reste limitée — une conséquence assumée et documentée de la taille de l'échantillon.
